\section*{Capítulo \ref{ch:b1:clastic-sediments} --- Sedimentos clásticos e seus ambientes}

\begin{solution}{exo:b1:clastic-sediments:1}
$\varphi = -\log_2 d$: \qty{0.18}{mm} dá $2.47$, uma areia fina (2 a 3); \qty{0.75}{mm} dá $0.42$, uma areia grossa (0 a 1); \qty{3}{mm} dá $-1.58$, um grânulo ($-1$ a $-2$); \qty{0.012}{mm} dá $6.38$, um silte fino (6 a 7); \qty{150}{mm} dá $-7.23$, um calhau ($-6$ a $-8$).
\end{solution}

\begin{solution}{exo:b1:clastic-sediments:2}
$M_z = (1.05 + 1.71 + 2.43)/3 = 1.73\,\varphi$, isto é, $2^{-1.73} =
\qty{0.30}{mm}$, uma areia média. $\sigma_I = (2.43 - 1.05)/4 + (2.98 - 0.62)/6.6
= 0.345 + 0.358 = 0.70\,\varphi$: moderadamente bem selecionada (0.50 a 0.71), logo abaixo do limite de moderadamente selecionada.
\end{solution}

\begin{solution}{exo:b1:clastic-sediments:3}
$q_b = (1 - 0.4) \times 0.02 \times (\num{0.01}/60)/2 = \qty{1.0e-6}{m^2.s^{-1}}$
de grãos sólidos. Em um dia, $\num{1.0e-6} \times 86400 = \qty{0.0864}{m^3}$ por
metro de crista, ou $0.0864 \times 2650 = \qty{229}{kg}$ de quartzo.
\end{solution}

\begin{solution}{exo:b1:clastic-sediments:4}
(a)~50 nós são $50 \times 1852/3600 = \qty{25.7}{m.s^{-1}}$; 12 nós, \qty{6.17}{m.s^{-1}}. (b)~$\qty{100000}{m}/\qty{6.17}{m.s^{-1}} =
\qty{16200}{s}$, cerca de \qty{4.5}{h}. (c)~480 milhas são \qty{772}{km}, 400 milhas
\qty{644}{km}.
\end{solution}

\begin{solution}{exo:b1:clastic-sediments:5}
(a)~$(\varphi_{84} - \varphi_{16})/2 = 2 \times 0.9945\,\sigma/2 =
0.9945\,\sigma$, e $(\varphi_{95} - \varphi_5)/3.3 = 2 \times
1.6449\,\sigma/3.3 = 0.9969\,\sigma$: ambos estimam $\sigma$ dentro de um terço de um por cento. (b)~Um sedimento real não precisa ser normal em unidades phi: uma pequena mistura de grãos muito mais grossos ou muito mais finos muda muito $\varphi_5$ ou $\varphi_{95}$, e pouco $\varphi_{16}$ e $\varphi_{84}$, de modo que a segunda estimativa cresce enquanto a primeira não. As duas diferem mais para misturas, como uma areia com algum cascalho e lama, ou um depósito bimodal.
\end{solution}

\begin{solution}{exo:b1:clastic-sediments:6}
(a)~$\sin\theta = T/L = 1.5/15 = 0.1$, logo $\theta = \ang{5.74}$. (b)~$a =
c\tan\theta = 0.5 \times 0.1005 = \qty{0.050}{cm}$ por minuto, ou
\qty{3.0}{cm} por hora. (c)~Um leito elevado em centímetros por hora recebe areia que cai da suspensão muito depressa: uma corrente muito carregada que desacelera, como uma enchente em declínio ou a cauda de uma \omterm{def:b1:clastic-sediments:turbidite}{corrente de turbidez}.
\end{solution}

\begin{solution}{exo:b1:clastic-sediments:7}
(a)~$(350 + 20 + 5 + 340 + 15 + 30)/6 = 760/6 = 126.7$: uma \omterm{def:b1:geological-maps:attitude}{direção} de sudeste, embora cada leitura esteja a menos de \ang{20} do norte. (b)~$S =
\sum\sin\alpha_i = 0.672$, $C = \sum\cos\alpha_i = 5.692$, logo $\bar\alpha =
\operatorname{atan2}(0.672, 5.692) = \ang{6.7}$, e $\bar R =
\sqrt{0.672^2 + 5.692^2}/6 = 5.732/6 = 0.955$. (c)~A média aritmética trata 350 e 005 como distantes, enquanto no círculo eles estão a \ang{15} um do outro: \omterm{def:b1:spherical-astronomy:horizontal}{azimutes} devem ser médios como vetores unitários. O $\bar R$ grande mostra que as leituras se agrupam firmemente em torno de \ang{7}.
\end{solution}

\begin{solution}{exo:b1:clastic-sediments:8}
A tem mais de \qty{5}{\percent} de argila: imaturo, e também imaturo na composição, com seu \omterm{def:b1:rock-forming-minerals:groups}{feldspato} e fragmentos de rocha; depositado depressa perto de sua fonte, em um leque aluvial ou uma planície de inundação, ou por uma \omterm{def:b1:clastic-sediments:turbidite}{corrente de turbidez}. B tem menos de \qty{5}{\percent} de argila e é bem selecionado ($\sigma_I < 0.5\,\varphi$), mas seus grãos são subangulosos ($\rho < 3$): maduro, por exemplo uma barra de rio ou uma areia de plataforma retrabalhada por correntes. C não tem argila, é muito bem selecionado e seus grãos são arredondados ($\rho > 3$): supermaduro, e com \qty{98}{\percent} de quartzo também maduro na composição, como numa praia ou num campo de \omterm{def:b1:clastic-sediments:bedform}{dunas} de deserto.
\end{solution}

\begin{solution}{exo:b1:clastic-sediments:9}
(a)~$1.5/\sqrt{9.81 \times 0.3} = 0.87$; $1.0/\sqrt{9.81 \times 2} = 0.23$.
(b)~$h = U^2/g = 1.5^2/9.81 = \qty{0.23}{m}$; $U = \sqrt{9.81 \times 2} =
\qty{4.43}{m.s^{-1}}$. (c)~O número de Froude cresce à medida que a profundidade diminui para uma dada velocidade: lâminas de água finas e rápidas, como o refluxo de uma onda numa praia, atingem facilmente números de Froude perto de 1, enquanto um rio de vários metros de profundidade precisaria de vários metros por segundo, que poucos rios atingem exceto em enchentes.
\end{solution}

\begin{solution}{exo:b1:clastic-sediments:10}
(a)~Com os pontos médios $-0.5, 0.5, \ldots, 4.5$: $m = (-1 + 4 + 45 + 100 + 52.5
+ 22.5)/100 = 2.23\,\varphi$; os desvios quadrados ponderados pelas porcentagens somam \num{107.71}, logo $s = \sqrt{1.0771} = 1.04\,\varphi$: mal selecionada. (b)~A porcentagem acumulada mais grossa é 0, 2, 10, 40, 80, 95 e
100 em $-1, 0, \ldots, 5\,\varphi$. Então $\varphi_5 = 0 + 3/8 = 0.375$,
$\varphi_{16} = 1 + 6/30 = 1.20$, $\varphi_{50} = 2 + 10/40 = 2.25$,
$\varphi_{84} = 3 + 4/15 = 3.27$, $\varphi_{95} = 4.0$. $M_z = (1.20 + 2.25 +
3.27)/3 = 2.24\,\varphi$ e $\sigma_I = (3.27 - 1.20)/4 + (4.0 - 0.375)/6.6 =
0.517 + 0.549 = 1.07\,\varphi$. (c)~Os dois métodos concordam em cerca de
\num{0.01} na média e \num{0.03} na \omterm{def:b1:clastic-sediments:sorting}{seleção}: uma areia fina ($2.24\,\varphi$ equivale a \qty{0.21}{mm}), mal selecionada.
\end{solution}

\begin{solution}{exo:b1:clastic-sediments:11}
(a)~$U = \sqrt{0.02 \times 9.81 \times 50 \times 0.005/0.003} =
\qty{4.04}{m.s^{-1}}$, ou 7.9 nós. (b)~$(\Delta\rho/\rho)\,h = U^2
C_f/(gS) = 25.7^2 \times 0.003/(9.81 \times 0.005) = \qty{40}{m}$: por exemplo,
uma corrente de \qty{400}{m} de espessura com excesso de densidade de \qty{10}{\percent}, o que mostra quão grande e densa deve ter sido a corrente de 1929. (c)~A velocidade varia com $\sqrt S$: $4.04/\sqrt 5 = \qty{1.81}{m.s^{-1}}$. A corrente mais lenta tem turbulência mais fraca e deposita seus grãos mais grossos; torna-se menos densa, desacelera ainda mais e espalha sua carga pela planície como uma \omterm{def:b1:clastic-sediments:turbidite}{camada gradada}.
\end{solution}

\begin{solution}{exo:b1:clastic-sediments:12}
(a)~$0.7 \times \qty{5e10}{kg}/\qty{2000}{kg.m^{-3}} = \qty{1.75e7}{m^3}$ por ano. (b)~Espalhado por uma frente de \qty{30}{m} de altura e \qty{100}{km} de largura, avança a frente em $\num{1.75e7}/(30 \times \num{1e5}) = \qty{5.8}{m}$ por ano. (c)~$a = v\tan\alpha = 5.83 \times \tan\ang{0.5} = \qty{0.051}{m}$, cerca de \qty{5}{cm} por ano, cem vezes menos que o avanço; a sucessão de \qty{30}{m} forma-se em $30/0.051 = 590$ anos.
\end{solution}

\begin{solution}{pb:b1:clastic-sediments:1}
\textbf{1.} Unidade~1: $(5.0 + 6.6 + 8.4)/3 = 6.67\,\varphi$, \qty{0.0098}{mm}
(\qty{9.8}{\micro m}); unidade~2: $3.47\,\varphi$, \qty{0.090}{mm}; unidade~3:
$1.80\,\varphi$, \qty{0.29}{mm}; unidade~4: $0.73\,\varphi$, \qty{0.60}{mm}.

\textbf{2.} Unidade~1: $(8.4 - 5.0)/4 + (9.6 - 4.2)/6.6 = 0.850 + 0.818 =
1.67\,\varphi$; unidade~2: $0.550 + 0.545 = 1.10$; unidade~3: $0.275 + 0.273 = 0.55$;
unidade~4: $0.725 + 0.712 = 1.44$.

\textbf{3.} Unidade~1: silte fino, mal selecionado; unidade~2: areia muito fina, mal
selecionada; unidade~3: areia média, moderadamente bem selecionada; unidade~4: areia
grossa, mal selecionada.

\textbf{4.} Unidade~3. Na foz do rio, a corrente e as ondas movem a areia repetidamente
e levam embora os finos, enquanto a lama do pró-delta se deposita por suspensão como uma
mistura de todas as granulometrias finas que lá chegam.

\textbf{5.} O tamanho $8\,\varphi$ fica $1.4/1.8$ do caminho de $\varphi_{50}$ a
$\varphi_{84}$, então $50 + (1.4/1.8) \times 34 = 76$ por cento da amostra é mais grossa
e cerca de \qty{24}{\percent} é de tamanho argila. Partículas tão finas só se depositam
em água quase parada: o local ficava abaixo do alcance de ondas e correntes.

\textbf{6.} Um canal transporta tanto cascalho, rolado no leito durante as cheias, quanto
areia; o depósito mistura o resíduo de cascalho do fluxo mais profundo com a areia dos
fluxos mais baixos, por isso cobre muitas classes.

\textbf{7.} A média cresce de \qty{0.0098}{mm} para \qty{0.090}{mm} e \qty{0.29}{mm}: a
sucessão engrossa para cima.

\textbf{8.} Unidade~1, lama laminada: o pró-delta; unidade~2, siltes e areias alternados:
a frente deltaica inferior; unidade~3, areia com \omterm{def:b1:clastic-sediments:cross-bedding}{estratificação cruzada}: as barras de
desembocadura da frente deltaica superior.

\textbf{9.} Cada zona passou a repousar sobre a que antes ficava à sua frente: o delta
avançou mar adentro, por cima do local.

\textbf{10.} Pelo menos \qty{21}{m}, a espessura das unidades 1 a 3, já que a sucessão
preencheu a coluna de água, desde o fundo do pró-delta até a barra de desembocadura perto da
superfície; mais, pois as lamas foram compactadas desde então
(\cref{ch:b1:diagenesis-basins}).

\textbf{11.} Um canal distributário da planície deltaica, cortando a barra de
desembocadura à medida que o delta avançava; seus seixos são sua \omterm{def:b1:erosion-transport:load}{carga de fundo}.

\textbf{12.} Não: uma barra de meandro começa com uma superfície de erosão e um resíduo e
afina para cima, até a lama de planície de inundação. A frente deltaica engrossa para
cima a partir da lama.

\textbf{13.} $\qty{16}{m}/\num{4000}\text{ anos} = \qty{4.0}{mm}$ por ano.

\textbf{14.} $v = a/\tan\alpha = 0.0040/\tan\ang{0.5} = 0.0040/0.008727 =
\qty{0.458}{m}$ por ano.

\textbf{15.} $0.458 \times 4000 = \qty{1830}{m}$, cerca de \qty{1.8}{km}.

\textbf{16.} $\qty{21}{m}/\qty{4.0}{mm.yr^{-1}} = 5250$ anos.

\textbf{17.} A frente deltaica recebe a maior parte da areia perto das desembocaduras, o
pró-delta muito menos e a planície deltaica pouco, exceto nas cheias; as desembocaduras
ativas também se deslocam de um lugar a outro, então um ponto pode receber muito por um
tempo e depois quase nada. Uma taxa única entre duas camadas datadas é uma média sobre
essas mudanças, e as lamas perderam parte de sua espessura por compactação.

\textbf{18.} A soma é 1544, então a média aritmética é \ang{128.7}, uma \omterm{def:b1:geological-maps:attitude}{direção} sudeste:
sem sentido, já que as leituras se espalham em torno do norte e os valores perto de 340
e perto de 010 são médios como se estivessem distantes.

\textbf{19.} $S = \sum\sin\alpha_i = 1.698$ e $C = \sum\cos\alpha_i =
11.248$.

\textbf{20.} $\bar R = \sqrt{1.698^2 + 11.248^2}/12 = 11.376/12 = 0.948$: as leituras se
agrupam firmemente em torno de uma só \omterm{def:b1:geological-maps:attitude}{direção}.

\textbf{21.} Os rios fluíam, e o delta avançava, em direção a aproximadamente norte
(\cref{prop:b1:clastic-sediments:azimuths}); a terra ficava ao sul.

\textbf{22.} Porque as estratificações frontais foram medidas em camadas que já não são
horizontais: suas direções de \omterm{def:b1:geological-maps:attitude}{mergulho} são as das camadas inclinadas mais a inclinação.
Girar cada estratificação frontal em torno do \omterm{def:b1:field-lab-methods-1:trend-plunge}{rumo} do acamamento pelo \omterm{def:b1:geological-maps:attitude}{mergulho} do
acamamento restaura o acamamento e, com ele, as estratificações frontais à sua \omterm{def:b1:geological-maps:attitude}{atitude}
original.

\textbf{23.} Correntes de maré entram e saem, e constroem séries que mergulham em duas
direções opostas, como na \omterm{def:b1:clastic-sediments:cross-bedding}{estratificação cruzada} em espinha de peixe; o diagrama de
rosácea teria dois lobos cerca de \ang{180} afastados. As doze leituras formam um só
lobo, com $\bar R$ perto de 1: uma corrente de rio fluindo em um sentido.

\textbf{24.} $\bar\alpha = \operatorname{atan2}(1.698, 11.248) = \ang{8.6}$: o \omterm{def:b1:spherical-astronomy:horizontal}{azimute}
médio da \omterm{def:b1:clastic-sediments:cross-bedding}{paleocorrente} da unidade~3 é \ang{9}, um pouco a leste do norte.
\end{solution}
